\section{Distributed Training \& Communication}

\smallskip\noindent\textit{Study Guide companion: Study Guide Section 12 (Hardware Deep Dive) covers the interview-ready version of this material.}
\smallskip

\begin{whythis}
A single GPU at 90\% utilization means nothing if 1000 GPUs spend 40\% of their time on communication.
At frontier scale---millions of accelerators across TPU, Trainium, and GPU clusters---the relevant metric
is not per-device FLOP/s but cluster-level intelligence-per-Watt.
Every wasted cycle waiting for a gradient synchronization or an expert routing token is a wasted Joule
multiplied by a million.
Distributed systems knowledge is therefore not a peripheral skill for a research engineer at this scale;
it is the primary lever for translating compute budget into capability gains.
RLVR training compounds this: unlike pretraining, it interleaves inference and training, produces
variable-length outputs that create load imbalance, and---for MoE architectures---relies on all-to-all
communication patterns that scale far worse than the all-reduce used in dense pretraining.
Understanding where the bubbles are, and which ones are worth eliminating, is the job.
\end{whythis}

% ----------------------------------------------------------------
\subsection*{Conceptual Foundation}
% ----------------------------------------------------------------

\begin{thinkfirst}{Parallelism Strategies: Memory Math for a 70B Model}
Before reading further, work through the following from first principles.
You have a 70B-parameter model (assume all weights in \texttt{bfloat16}) and 8 GPUs each with 80 GB HBM.

\medskip
\textbf{Data Parallel (DDP).}
Every GPU holds a full copy of the model.
70B parameters $\times$ 2 bytes/param = 140 GB per GPU.
Each GPU also holds gradients (140 GB) and, with Adam, two optimizer state tensors (280 GB).
Total per GPU: $\sim$560 GB.
\emph{Does it fit?} No --- 560 GB $\gg$ 80 GB.
DDP alone is ruled out for a 70B model on 80 GB GPUs.

\medskip
\textbf{Tensor Parallelism (TP).}
Split each weight matrix along one dimension across $T$ GPUs.
An attention projection of shape $[d, d_{\text{head}} \cdot h]$ is split column-wise; each GPU holds
$1/T$ of the columns.
With $T = 8$, parameter memory per GPU: $140/8 = 17.5$ GB --- fits comfortably.
Cost: each forward pass requires an \texttt{all-reduce} after every tensor-parallel layer
(to combine partial sums).
This is a \emph{synchronous, latency-sensitive} communication on the critical path.

\medskip
\textbf{Pipeline Parallelism (PP).}
Assign contiguous layer groups to different GPUs (stages).
With 8 stages and 80 layers, each GPU holds 10 layers $\approx 17.5$ GB --- fits.
Cost: \emph{pipeline bubbles}.
With $M$ micro-batches and $P$ stages, the bubble fraction is $(P-1)/(M+P-1)$.
For $P=8$ and $M=8$: bubble fraction $= 7/15 \approx 47\%$ --- nearly half the step is idle.
Increasing $M$ reduces the bubble but increases memory for activations.

\medskip
\textbf{Fully Sharded Data Parallel (FSDP).}
Shard \emph{all} tensors (parameters, gradients, optimizer states) across all $N$ GPUs.
Memory per GPU: $\text{total} / N$.
For 70B + Adam: $\sim$560 GB / 8 = 70 GB --- barely fits with careful activation checkpointing.
Protocol: (1) \texttt{all-gather} the full parameter tensor for each layer before its forward pass,
(2) compute forward, (3) compute backward using the gathered weights,
(4) \texttt{reduce-scatter} gradients so each GPU receives its shard's gradient,
(5) update local shard.
The gather and reduce-scatter double the communication volume relative to DDP's all-reduce, but
overlap with computation when implemented carefully.

\medskip
\textbf{Hybrid strategies.}
Production systems combine all three.
A typical recipe: TP within a node (NVLink), PP across nodes (InfiniBand), FSDP across all replicas.
The reason TP stays intra-node is that tensor-parallel all-reduces are on the critical path of every
forward pass --- they must complete before the next layer starts.
NVLink bandwidth ($\sim$900 GB/s bidirectional on H100 NVLink4) can absorb this; InfiniBand cannot.

\medskip
\textbf{Design question:} For the 70B model on 8$\times$80 GB GPUs, sketch a strategy that fits in
memory and minimizes compute bubbles.
What does your communication graph look like?
Where is the bottleneck?
\end{thinkfirst}

\begin{thinkfirst}{Why RLVR Distribution Is Harder Than Pretraining}
RLVR training is not a single workload --- it is a pipeline of at least three distinct
compute phases that must be orchestrated across the same (or overlapping) hardware.
Before reading the analysis, ask yourself: which phase is memory-bandwidth-bound and which is
compute-bound?
Why does this distinction matter for how you schedule work across GPUs?

\medskip
\textbf{Phase 1 --- Rollout generation (inference).}
The policy generates $G$ completions per prompt (GRPO uses $G \sim 8$--$16$).
Generation is autoregressive: each token step requires loading the full KV cache plus the model
weights, performing a small matrix-vector multiply, and sampling.
This is \emph{memory-bandwidth-bound}, not compute-bound.
A single A100 has $\sim$2 TB/s HBM bandwidth; with a 7B model at bfloat16, loading weights takes
$\sim$7 ns per byte but the flop count per token is tiny relative to that bandwidth.
GPU utilization during generation: typically 30--60\%, even with batching.

\medskip
\textbf{Phase 2 --- Reward evaluation.}
Each completion must be scored by a reward model (or verified by a verifier in RLVR math settings).
This is a separate forward pass.
If the reward model differs in architecture from the policy, it may need its own GPU partition or
must time-share.
Scheduling these evaluations without stalling the generation pipeline requires careful batching.

\medskip
\textbf{Phase 3 --- Policy update (training).}
Once all $G$ completions for a batch of prompts are collected, advantages are computed and a gradient
step is taken.
This is \emph{compute-bound}: large matrix multiplications in the backward pass.
GPU utilization should be high here --- but only if the batch is large enough and
fully padded, which leads to the next problem.

\medskip
\textbf{Variable-length outputs and load imbalance.}
Completions vary in length.
A prompt that terminates early leaves its GPU idle while others finish.
The training step cannot begin until \emph{all} $G$ completions per prompt are done
(needed for advantage normalization in GRPO).
This creates a synchronization barrier at the end of every generation round.
The slowest completion in the batch determines when training can start --- the classic straggler problem.

\medskip
\textbf{MoE expert parallelism under RLVR.}
Dense models use all-reduce: every device sends/receives the same aggregated gradient.
MoE models route each token to $k$ of $E$ experts; experts are sharded across GPUs.
This requires \texttt{all-to-all} communication: device $i$ sends a \emph{different} chunk to
every other device.
During pretraining, routing distributions are approximately stationary.
During RLVR, the policy changes at every step, and routing patterns shift dynamically ---
some experts become more popular, others are underutilized.
This creates \emph{dynamic} load imbalance on top of the static straggler problem.

\medskip
\textbf{Design question:} Sketch a timeline diagram for one RLVR step on 4 GPUs using GRPO with
$G=8$ and a 7B MoE policy.
Label each phase and mark the synchronization barriers.
Where are the bubbles?
Which ones can be overlapped and which cannot?
\end{thinkfirst}

\begin{thinkfirst}{Communication Patterns: All-Reduce vs.\ All-to-All}
Before reading the derivations, estimate: for 8 GPUs exchanging 1 GB of gradient data,
how long does a ring all-reduce take if each GPU-to-GPU link has 100 GB/s bandwidth?
Then read below and check your estimate.

\medskip
\textbf{Ring all-reduce.}
Arrange $N$ devices in a ring.
Each device starts with a tensor of size $S$ bytes.
The ring all-reduce proceeds in two phases, each with $N-1$ steps.

\emph{Reduce-scatter phase:} Each device sends a chunk of size $S/N$ to the next device in the ring,
receives a chunk from the previous device, and adds it to its local buffer.
After $N-1$ steps, each device holds one fully reduced chunk of the global sum.

\emph{All-gather phase:} Each device sends its reduced chunk to the next device in the ring.
After $N-1$ steps, every device holds all $N$ reduced chunks, i.e., the full sum.

Total data sent per device: $2 \cdot (N-1)/N \cdot S$ bytes.
As $N \to \infty$, this approaches $2S$ bytes --- essentially twice the tensor size, regardless of $N$.
Bandwidth utilization: $(N-1)/N$.
For $N=8$, utilization $= 7/8 = 87.5\%$.
Latency: $2(N-1) \cdot \alpha + 2(N-1)/N \cdot S / \beta$,
where $\alpha$ is per-hop latency and $\beta$ is per-link bandwidth.
Ring all-reduce is bandwidth-optimal but has $O(N)$ latency --- poor for large $N$ at small message sizes.

\medskip
\textbf{Tree (recursive halving/doubling) all-reduce.}
Devices are paired, then groups of 4, then 8, etc.
Number of steps: $\log_2 N$.
Latency: $O(\log N \cdot \alpha)$ --- much better for small messages.
But each step doubles the message size being transmitted.
Bandwidth utilization is lower because early steps transmit small messages inefficiently.
Tree all-reduce wins when $\alpha$ dominates (small tensors, high-latency links).
Ring all-reduce wins when $\beta$ dominates (large tensors, high-bandwidth links like NVLink).

\medskip
\textbf{All-to-all (MoE routing).}
Every device sends a \emph{different} chunk to every other device.
Total data sent per device: $(N-1) \cdot c$ bytes, where $c$ is the chunk size per destination.
Total cluster communication: $N \cdot (N-1) \cdot c$ bytes --- $O(N^2)$.
Compare to ring all-reduce: $O(N \cdot S)$ bytes cluster-total, which for uniform sharding is
$O(S)$ per device regardless of $N$.
All-to-all scales catastrophically with $N$ for fixed per-device capacity.
For MoE with $E$ experts sharded across $N$ GPUs, each forward pass requires two all-to-all calls
(dispatch tokens to experts, receive results back).
Under RLVR, routing is non-uniform and changes every step --- some all-to-all messages are large,
some are tiny, and the imbalance stresses the network fabric in unpredictable ways.

\medskip
\textbf{Hardware implications.}
NVIDIA NVLink provides all-to-all bandwidth within a node that is roughly symmetric.
AMD Infinity Fabric (used with RCCL, the ROCm equivalent of NCCL) has a different topology:
bandwidth between two specific GPUs depends on which Infinity Fabric links connect them,
which may be asymmetric.
This matters for all-to-all: if expert 3 is always sending to GPU 7 via a slower link,
that link becomes a bottleneck every step.
RCCL's all-to-all implementation may need topology-aware routing that NCCL does not,
because NVLink's switched topology is more uniform.
This is a concrete hardware-software co-design question relevant to AMD evaluation (Chapter 1).

\medskip
\textbf{Design question:} For a 16-GPU cluster running MoE with 64 experts, compute the all-to-all
communication volume per step if each token's expert routing vector is 1 KB and you process 4096
tokens per step.
Compare this to the gradient all-reduce volume for a 7B dense model.
At what batch size does the all-to-all dominate?
\end{thinkfirst}

% ----------------------------------------------------------------
\subsection*{Exercises}
% ----------------------------------------------------------------

\begin{exercise}{Ring All-Reduce from Scratch (3--4 hours)}
Implement ring all-reduce using raw Python multiprocessing on a single machine.
Do not use \texttt{torch.distributed.all\_reduce} for the implementation itself --- use it only as a
reference for correctness checking.

\medskip
\textbf{Setup.}
Spawn 4 processes using \texttt{torch.multiprocessing} or Python's \texttt{multiprocessing}.
Each process represents one ``device'' and holds a local tensor of size $M$ bytes.
Use shared memory or Unix sockets for inter-process communication.

\medskip
\textbf{Step 1 --- Reduce-scatter.}
Divide each tensor into 4 chunks.
Implement the ring: process $i$ sends chunk $i$ to process $(i+1) \bmod 4$,
receives from process $(i-1) \bmod 4$, and accumulates.
Repeat for $N-1 = 3$ steps, rotating which chunk is sent.
After this phase, process $i$ holds the globally summed version of chunk $i$.

\medskip
\textbf{Step 2 --- All-gather.}
Process $i$ sends its reduced chunk to process $(i+1) \bmod 4$.
After $N-1$ steps, every process holds the full summed tensor.
Verify correctness against \texttt{torch.distributed.all\_reduce}.

\medskip
\textbf{Step 3 --- Benchmarking.}
Measure achieved bandwidth vs.\ theoretical for message sizes:
1 KB, 10 KB, 100 KB, 1 MB, 10 MB, 100 MB, 1 GB.
Plot bandwidth (GB/s) vs.\ message size (log scale).
You should see two distinct regimes:
\begin{itemize}
  \item \emph{Latency-dominated} (small messages): bandwidth is far below peak because
        setup and synchronization overhead dominate.
  \item \emph{Bandwidth-dominated} (large messages): throughput approaches the
        theoretical limit of your inter-process transport.
\end{itemize}
Identify the crossover point in bytes.
Compute theoretical bandwidth for your transport and compare.

\medskip
\textbf{Step 4 --- NCCL comparison.}
Repeat the bandwidth plot using \texttt{torch.distributed.all\_reduce} with NCCL backend on actual
GPUs (or \texttt{gloo} backend if only CPU is available).
At what message size does NCCL reach $>$80\% of the theoretical NVLink bandwidth?
How does NCCL's crossover point compare to your implementation's?

\medskip
\textbf{Deliverables.} Bandwidth-vs-message-size plot with both curves.
Annotated crossover point.
One paragraph explaining the shape of each curve in terms of $\alpha$--$\beta$ cost model.
\end{exercise}

\begin{exercise}{NCCL Profiling on Multi-GPU Instance (2--3 hours)}
Rent a multi-GPU instance (2$\times$H100 $\approx$ \$5/hr, or 4$\times$A100 $\approx$ \$6/hr).
Budget approximately \$15--20 for this exercise.

\medskip
\textbf{Benchmark setup.}
Use \texttt{torch.distributed} with NCCL backend.
Write a script that launches all-reduce for message sizes:
1 MB, 10 MB, 100 MB, 1 GB.
For each size: run 10 warmup iterations, then 50 timed iterations.
Record mean and standard deviation of latency.
Compute achieved bandwidth as $2(N-1)/N \cdot S / \text{latency}$.

\medskip
\textbf{Profiling tools.}
Enable \texttt{NCCL\_DEBUG=INFO} and redirect stderr to a log file.
Parse the log to find: which NCCL algorithm is selected for each message size (ring, tree, or
collnet), and what the reported bus bandwidth is.
Optionally, run under \texttt{nsys profile} to get a GPU timeline.
In the timeline, identify: NCCL kernel launches, data copy operations, and synchronization points.

\medskip
\textbf{Overhead decomposition.}
For the 1 MB case, what fraction of the all-reduce time is:
\begin{enumerate}
  \item Kernel launch and setup overhead?
  \item Actual data transfer?
  \item Synchronization/barrier?
\end{enumerate}
How does this fraction change at 1 GB?
At what message size does NCCL switch algorithms?
Does the switch correspond to the crossover you found in the previous exercise?

\medskip
\textbf{Scaling.}
If your instance has 4 GPUs, run all-reduce on 2, 3, and 4 GPUs for a fixed 100 MB tensor.
Does bandwidth scale as $(N-1)/N$ predicts?
If not, why not?

\medskip
\textbf{Deliverables.} Bandwidth-vs-message-size table with algorithm annotations.
Overhead decomposition bar chart for 1 MB and 1 GB.
Written answer: at what message size does NCCL reach $>$80\% of peak NVLink bandwidth?
\end{exercise}

\begin{exercise}{Toy FSDP: Manual Sharding of GPT-2 Small (4--6 hours)}
Implement FSDP manually for GPT-2 small (117M parameters) on 2 GPUs.
Do not use PyTorch's built-in \texttt{FullyShardedDataParallel} for the core implementation;
use it only for comparison.

\medskip
\textbf{Step 1 --- Shard parameters.}
For each parameter tensor $\theta \in \mathbb{R}^{d}$, GPU $i$ stores
$\theta_i = \theta[i \cdot d/2 : (i+1) \cdot d/2]$.
Do this for all parameters: weights, biases, embedding tables.
Verify that memory per GPU is approximately half of full-model memory.

\medskip
\textbf{Step 2 --- All-gather before forward.}
Before each layer's forward pass, issue \texttt{dist.all\_gather} to reconstruct the full parameter
tensor on both GPUs.
Perform the forward computation.
Immediately free the gathered tensor to reclaim memory (critical: if you keep all gathered tensors,
you have replicated the full model and defeated the purpose).

\medskip
\textbf{Step 3 --- Backward and reduce-scatter.}
PyTorch's autograd will compute gradients with respect to the gathered tensor.
After the backward pass, issue \texttt{dist.reduce\_scatter} to distribute gradient shards:
GPU $i$ accumulates the gradient for its parameter shard.
Each GPU now has the gradient for its shard only.

\medskip
\textbf{Step 4 --- Local optimizer update.}
Each GPU runs its optimizer (Adam) on its local parameter shard using its local gradient shard.
Optimizer state (momentum, variance) is also sharded and never needs to be gathered.
This is the memory saving: optimizer state is $2\times$ parameter memory for Adam,
and it is always sharded.

\medskip
\textbf{Step 5 --- Measurement.}
Compare peak memory per GPU for: (a) full replication (DDP), (b) your manual FSDP,
(c) PyTorch FSDP.
Measure throughput (samples/sec) for a fixed batch.
Where is your implementation slower than PyTorch FSDP?
Profile to find: is the bottleneck in the all-gather, the forward pass, or the reduce-scatter?

\medskip
\textbf{Deliverables.}
Memory comparison table.
Throughput comparison.
Profiler output annotated with where time is spent relative to PyTorch FSDP.
One paragraph: what does PyTorch FSDP do that your implementation does not, based on the profiler evidence?
\end{exercise}

\begin{exercise}{RLVR Distribution: Profiling Generation Bubbles (4--6 hours)}
Take a minimal RLVR training loop (nanoRLVR from Chapter 8, or build a stripped-down GRPO loop
around a 1B--3B model) and instrument it for distribution profiling on 2 GPUs.

\medskip
\textbf{Setup.}
Split rollout generation across 2 GPUs (each GPU generates $G/2$ completions per prompt).
Both GPUs participate in the training step.
Use a dataset of math problems (e.g., GSM8K) with verifier-based rewards so reward evaluation is
deterministic and fast.

\medskip
\textbf{Instrumentation.}
Wrap each phase with \texttt{torch.profiler} or manual \texttt{time.perf\_counter} checkpoints:
\begin{itemize}
  \item $t_{\text{gen\_start}}$, $t_{\text{gen\_end}}$: generation phase per GPU
  \item $t_{\text{reward\_start}}$, $t_{\text{reward\_end}}$: reward evaluation
  \item $t_{\text{sync}}$: time waiting at the synchronization barrier (straggler cost)
  \item $t_{\text{train\_start}}$, $t_{\text{train\_end}}$: backward pass
  \item $t_{\text{comm}}$: gradient communication time
\end{itemize}

\medskip
\textbf{Bubble identification.}
For a batch of 16 prompts with $G=8$:
compute the fraction of total step time spent in each phase.
The useful compute fraction is $t_{\text{train}} / t_{\text{step}}$.
Everything else is overhead: generation (unavoidable but schedulable), reward (schedulable),
sync wait (the bubble), and gradient communication (reducible with overlap).
What fraction of the step is the synchronization wait (straggler bubble)?

\medskip
\textbf{Optimization.}
Propose and implement one optimization targeting the largest bubble you found.
Options:
\begin{itemize}
  \item \emph{Async generation:} start the next round of generation while the training backward pass
        runs on completed rollouts.
  \item \emph{Staggered scheduling:} assign shorter prompts to one GPU and longer prompts to the
        other to balance generation time.
  \item \emph{Early cutoff:} truncate the slowest completions once a timeout is reached
        (analyze the reward distribution impact).
\end{itemize}

\medskip
\textbf{Measurement.}
Before and after your optimization: report step time, useful compute fraction,
and synchronization wait fraction.
Target: reduce the synchronization bubble by $>$30\%.

\medskip
\textbf{Deliverables.}
Flame graph (torch profiler output) showing step phases and communication stalls.
Before/after utilization bar chart.
One paragraph: which bubble was largest, which optimization you chose, and why.
\end{exercise}

% ----------------------------------------------------------------
\subsection*{Intelligence-per-Watt at Cluster Scale}
% ----------------------------------------------------------------

For $N$ GPUs, define the useful compute fraction $\eta(N)$ as the fraction of aggregate FLOP/s
that produces gradients contributing to model improvement.
The rest is overhead: communication, synchronization barriers, and idle bubbles.

For \textbf{pretraining with all-reduce} (DDP or FSDP), $\eta(N)$ declines slowly with $N$
because ring all-reduce volume is $O(S)$ per device regardless of $N$.
The bottleneck is inter-node bandwidth; $\eta$ is typically $>$70\% even at thousands of GPUs
when network is not the bottleneck.

For \textbf{RLVR with GRPO synchronization}, $\eta(N)$ declines faster due to the straggler
barrier at the end of every generation round.
As $N$ grows, the probability that at least one GPU has a slow completion increases,
and the expected straggler delay grows as $O(\log N)$ for i.i.d.\ exponential completion times.
At large $N$, generation bubbles can consume $>$40\% of step time.

For \textbf{MoE with all-to-all routing}, $\eta(N)$ declines most steeply.
All-to-all volume scales as $O(N)$ per device (each device sends to all others), so communication
time grows linearly with $N$ while compute per device stays fixed.
There exists a critical $N^*$ beyond which adding more GPUs reduces $\eta$ faster than it adds
parallel throughput, making additional GPUs net-negative in terms of intelligence-per-Watt.

This is the scaling ceiling.
Identifying $N^*$ for each workload --- and building the measurement infrastructure to detect when
you are approaching it --- is a core research engineering responsibility at frontier scale.

% ----------------------------------------------------------------
\subsection*{Portfolio Project}
% ----------------------------------------------------------------

\begin{portfolio}{RLVR Training Distribution Profiler}
\textbf{The question.}
What fraction of RLVR training compute is wasted on distribution overhead, and which specific
overhead is most cost-effective to eliminate?

\medskip
\textbf{The monkey.}
Quantify the largest utilization gap in an RLVR training loop
(generation bubble, communication stall, GRPO synchronization wait, or MoE routing imbalance)
and implement one optimization that reduces it by a measurable amount.

\medskip
\textbf{The pedestal.}
A profiling harness that instruments a real RLVR training loop (minimum 1B-parameter policy,
real reward signal) and produces per-step breakdowns of:
generation time, reward evaluation time, synchronization wait time, gradient communication time,
and backward pass time.
The harness must produce flame graphs and utilization charts automatically.

\medskip
\textbf{Implementation plan.}
\begin{enumerate}
  \item Run the RLVR loop uninstrumented for 100 steps to establish baseline step time and throughput.
  \item Add \texttt{torch.profiler} instrumentation. Run 20 steps and export the profiler trace.
        Open the trace in Chrome \texttt{chrome://tracing} or TensorBoard profiler.
        Identify the top-3 time sinks per step.
  \item For the largest bubble ($>$10\% of step time), design one targeted optimization.
        Write a one-paragraph justification before implementing.
  \item Implement the optimization. Run 20 instrumented steps.
        Compare: step time, useful compute fraction ($\eta$), and peak GPU memory.
  \item Write the blog post.
        Required figures: (a) pre-optimization flame graph with bubble annotated,
        (b) post-optimization flame graph, (c) per-phase utilization bar chart before/after,
        (d) step time distribution histogram showing variance reduction from straggler mitigation.
\end{enumerate}

\medskip
\textbf{Verification criteria.}
\begin{enumerate}
  \item Flame graph shows at least one identifiable bubble $>$10\% of step time.
  \item Proposed optimization reduces that bubble by $>$30\% in measured wall time.
  \item Blog post is written for an ML engineer audience with sufficient implementation detail
        to reproduce the measurement setup and the optimization.
\end{enumerate}

\medskip
\textbf{Hardware budget.}
2$\times$H100 at $\approx$\$5/hr or 4$\times$A100 at $\approx$\$6/hr.
Profiling runs are short; budget $\approx$\$50 for this project within the \$200/month allocation.
\end{portfolio}

\newpage
